\section{Conclusions}\label{sec:conclusions}
The extended Zone architecture makes code--data identification an
explicit component of payload rewriting. Its symbolic and heuristic
path and its LLM-based alternative feed a shared deterministic
data-handling layer. This separation supports analysis of candidate
quality, reference handling, and downstream functionality as distinct
questions.

The transformation extension introduces a separate use of LLMs,
with zero-shot and two-shot conditions and explicit validation and
feedback. Assembly wrapping adds an implementation of an established
technique, while the x86 port remains development work until its
architecture-specific assessment is completed. These contributions
extend the conference design beyond its original transformation
catalogue and experimental baseline.

The assessment framework distinguishes construction, empirical
functionality, static verdicts, similarity, diversity, and cost.
Its purpose is to connect each conclusion to the evidence that can
support it. The historical measurements characterize the earlier
engine and corpus; completed extension results will characterize
the new paths under their stated configurations. Broader claims
about semantic equivalence, unseen inputs, or deployed endpoint
products require evidence outside those boundaries.
\draftnote{Before submission, add the principal verified findings and
the answers to RQ1--RQ4. Update this section and the abstract together;
remove development claims only when the corresponding assessment is
complete.}
